\section*{The 31st International Physics Olympiad}

\begin{center}
Translated and edited by Zheng Yongling, Jiang Zuimin and Lu Shenlong, Fudan University
\end{center}

\subsection*{Theoretical Problems}

\subsection*{Problem 1}

\textbf{Part A: The bungee jumper}

A bungee jumper is attached to one end of a long elastic rope, the other end of which is fixed to a high bridge. He falls from rest from the bridge towards the river below without touching the water. His mass is $m$, the natural length of the rope is $L$, the force constant of the rope (the force required to stretch it by $1\unitv{m}$) is $k$, and the strength of the gravitational field is $g$. Find expressions for the following quantities:

\begin{enumerate}
\item[(a)] the distance $y$ that the jumper falls before coming to instantaneous rest for the first time;
\item[(b)] the maximum speed $v$ that he attains during the fall;
\item[(c)] the time $t$ taken for the fall up to the first instantaneous rest.
\end{enumerate}

The jumper may be regarded as a particle attached to one end of the rope, and the mass of the rope is negligible compared with $m$. When the rope is stretched it obeys Hooke's law, and the resistance of the air may be neglected throughout the fall.

\textbf{Part B: The heat engine}

A heat engine operates between two bodies made of the same material. The temperatures of the two bodies are $T_A$ and $T_B$ ($T_A>T_B$), the mass of each body is $m$, and their specific heat capacities are constant and equal to $s$. The pressure of the two bodies is kept constant and no phase change occurs.

\begin{enumerate}
\item[(a)] Assuming that the engine extracts the maximum mechanical energy theoretically permitted from the system, find an expression for the final temperature $T$ reached by both bodies $A$ and $B$. Give the complete reasoning of the solution.
\item[(b)] Hence obtain an expression for the maximum work that can be obtained.
\item[(c)] Suppose that the heat engine operates between two tanks of water, each of volume $2.50\unitv{m^3}$, one at a temperature of $350\unitv{K}$ and the other at $300\unitv{K}$. Calculate the maximum mechanical energy that can be obtained.
\end{enumerate}

Given: specific heat capacity of water $=4.19\times10^3\unitv{J\cdot kg^{-1}\cdot K^{-1}}$, density of water $=1.00\times10^3\unitv{kg\cdot m^{-3}}$.

\textbf{Part C: Radioactivity and the age of the Earth}

Assume that when the Earth was formed the isotopes $^{238}\mathrm{U}$ and $^{235}\mathrm{U}$ were already present, but that their decay products were not. The decays of $^{238}\mathrm{U}$ and $^{235}\mathrm{U}$ are used to determine the age $T$ of the Earth.

\begin{enumerate}
\item[(a)] The isotope $^{238}\mathrm{U}$ decays with a half-life of $4.50\times10^{9}$ years. The half-lives of the remaining radioactive decay products in the decay chain are all much shorter, so that to a first approximation the presence of these products may be ignored; the chain terminates with the lead isotope $^{206}\mathrm{Pb}$. Using the half-life of $^{238}\mathrm{U}$ and the present number $^{238}N$ of $^{238}\mathrm{U}$ nuclei, express the number $^{206}n$ of $^{206}\mathrm{Pb}$ atoms produced by the radioactive decay (it is convenient to work in units of $10^{9}$ years).
\item[(b)] Similarly, $^{235}\mathrm{U}$ decays, after passing through a series of shorter-lived products, with a half-life of $0.710\times10^{9}$ years, terminating in the stable isotope $^{207}\mathrm{Pb}$. Write down the relation between $^{207}n$, $^{235}N$ and the half-life of $^{235}\mathrm{U}$.
\item[(c)] A mixed ore of lead and uranium is analysed with a mass spectrometer, and the relative concentrations of the lead isotopes $^{204}\mathrm{Pb}$, $^{206}\mathrm{Pb}$ and $^{207}\mathrm{Pb}$ in this ore are found to be in the ratio $1.00:29.6:22.6$. Since the isotope $^{204}\mathrm{Pb}$ is not radioactive, it can be used as a reference in the analysis. Analysis of a pure lead ore gives the relative concentrations of these three isotopes as $1.00:17.9:15.5$. Given that the ratio $^{238}N:^{235}N$ is $137:1$, derive a relation involving $T$.
\item[(d)] Assuming that the age $T$ of the Earth is much greater than the half-lives of both uranium isotopes, obtain an approximate value of $T$.
\item[(e)] Clearly the above approximation is not markedly greater than the longer of the half-lives of the uranium isotopes, but a more accurate value of $T$ can be obtained from this approximation. Hence estimate the age $T$ of the Earth to within $2\%$.
\end{enumerate}

\textbf{Part D: A spherical charge distribution}

In vacuum a charge $Q$ is distributed uniformly throughout a sphere of radius $R$.

\begin{enumerate}
\item[(a)] Derive the electric field strength at a distance $r$ from the centre of the sphere for the two cases $r\le R$ and $r>R$.
\item[(b)] Derive an expression for the total electrical energy associated with this charge distribution.
\end{enumerate}

\textbf{Part E: A rotating copper ring}

A circular ring made of fine copper wire rotates about its vertical diameter in the Earth's magnetic field. The magnetic flux density of the Earth's field at the position of the ring is $44.5\unitv{\mu T}$, and its direction makes an angle of $64\degree$ downwards with the horizontal. Given that the density of copper is $8.90\times10^3\unitv{kg\cdot m^{-3}}$ and its resistivity is $1.70\times10^{-8}\unitv{\Omega\cdot m}$, calculate the time required for its angular velocity to fall from its initial value to half of that value. Show your working. This time is much longer than the time of one revolution. Assume that friction due to the air and at the bearings is negligible, and neglect self-induction (although these effects ought not really to be neglected).

\subsection*{Solutions}

\subsection*{Problem 1}

\textbf{Part A: The bungee jumper}

(a) By conservation of energy,
\begin{equation*}
mgy=\frac{1}{2}k(y-L)^2 .
\end{equation*}
Solving,
\begin{equation*}
y=\frac{kL+mg+\sqrt{2mgkL+m^2g^2}}{k}
\end{equation*}
(the positive sign must be taken before the square root).

(b) The maximum speed occurs when the acceleration is zero, that is, when the forces balance. Let the extension of the rope at that instant be $x$; then
\begin{equation*}
mg=kx .
\end{equation*}
By the work--energy theorem,
\begin{equation*}
\frac{1}{2}mv_m^2=mg(L+x)-\frac{1}{2}kx^2 .
\end{equation*}
Solving these two equations,
\begin{equation*}
v_m=\sqrt{2gL+mg^2/k} .
\end{equation*}

(c) The time of the fall is the sum of the time $t_f$ of free fall and the time $t_s$ of simple harmonic motion of the elastic rope up to its maximum extension:
\begin{equation*}
t=t_f+t_s ,
\end{equation*}
\begin{equation*}
t_f=\sqrt{\frac{2L}{g}} .
\end{equation*}
The jumper enters the simple harmonic motion with the speed $v_t=gt_f=\sqrt{2gL}$ acquired in free fall, and his speed may be written as
\begin{equation*}
v=v_m\sin\omega t ,
\end{equation*}
where
\begin{equation*}
\omega=\frac{2\pi}{T}=\sqrt{\frac{k}{m}} ,
\end{equation*}
and
\begin{equation*}
v_t=v_m\sin\omega\tau ,
\end{equation*}
\begin{equation*}
\tau=\frac{1}{\omega}\arcsin\frac{v_t}{v_m}
=\frac{1}{\omega}\arcsin\frac{\sqrt{2gL}}{v_m} .
\end{equation*}
\begin{equation*}
t_s=\frac{T}{2}-\tau=\pi\sqrt{\frac{m}{k}}-\tau .
\end{equation*}
Hence
\begin{equation*}
t=t_f+t_s=\sqrt{\frac{2L}{g}}+\pi\sqrt{\frac{m}{k}}
-\frac{1}{\omega}\arcsin\frac{\sqrt{2gL}}{v_m} ,
\end{equation*}
\begin{equation*}
t=\sqrt{\frac{2L}{g}}+\sqrt{\frac{m}{k}}
\left[\pi-\arcsin\frac{\sqrt{2gL}}{\sqrt{2gL+mg^2/k}}\right] .
\end{equation*}
It may also be written in other equivalent forms.

\textbf{Part B: The heat engine}

In calculating the work obtainable, the friction losses in the operation of the engine are neglected.

Let $\Delta Q_1$ be the heat absorbed from body $A$, and $\Delta Q_2$ the heat released to body $B$; then
\begin{equation*}
\Delta Q_1=-ms\,\Delta T_1\qquad(\Delta T_1<0),
\end{equation*}
\begin{equation*}
\Delta Q_2=ms\,\Delta T_2\qquad(\Delta T_2>0).
\end{equation*}

(a) To obtain the maximum mechanical energy, the engine is taken to be a Carnot engine. By the second law of thermodynamics, throughout the working process,
\begin{equation*}
\frac{\Delta Q_1}{T_1}=\frac{\Delta Q_2}{T_2} ,
\end{equation*}
that is,
\begin{equation*}
-ms\int_{T_A}^{T_0}\frac{\dd T_1}{T_1}
=ms\int_{T_B}^{T_0}\frac{\dd T_2}{T_2} ,
\end{equation*}
\begin{equation*}
\ln\frac{T_A}{T_0}=\ln\frac{T_0}{T_B} .
\end{equation*}
Hence
\begin{equation*}
T_0=\sqrt{T_AT_B} .
\end{equation*}

(b)
\begin{align*}
W&=Q_1-Q_2=ms(T_A-T_0)-ms(T_0-T_B)\\
&=ms\left(T_A+T_B-2\sqrt{T_AT_B}\right),
\end{align*}
or
\begin{equation*}
W=ms\left(\sqrt{T_A}-\sqrt{T_B}\right)^2 .
\end{equation*}

(c) Numerical result:
\begin{align*}
W&=2.50\times1.00\times10^3\times4.19\times10^3
\left[350+300-2\sqrt{350\times300}\right]\unitv{J}\\
&=2.0\times10^7\unitv{J}=20\unitv{MJ} .
\end{align*}

\textbf{Part C: Radioactivity and the age of the Earth}

(a)
\begin{equation*}
N=N_0\,2^{-t/\tau} ,
\end{equation*}
where $N_0$ is the original number of atoms and $\tau$ is the half-life.
\begin{equation*}
n=N_0-N_0\,2^{-t/\tau}=N_0\left(1-2^{-t/\tau}\right).
\end{equation*}
Expressed in terms of the present number $N$ of atoms, this becomes
\begin{equation*}
n=N\left(2^{t/\tau}-1\right).
\end{equation*}
Therefore
\begin{equation*}
^{206}n={}^{238}N\left(2^{t/4.50}-1\right),
\end{equation*}
where $t$ is measured in units of $10^{9}$ years.

(b) In the same way,
\begin{equation*}
^{207}n={}^{235}N\left(2^{t/0.710}-1\right).
\end{equation*}

(c)
\begin{equation*}
\frac{^{206}n}{^{207}n}
=\frac{^{238}N}{^{235}N}\cdot
\frac{2^{t/4.50}-1}{2^{t/0.710}-1} ,
\end{equation*}
that is,
\begin{equation*}
\frac{29.6-17.9}{22.6-15.5}
=137\,\frac{2^{T/4.50}-1}{2^{T/0.710}-1} ,
\end{equation*}
or
\begin{equation*}
0.0120\left(2^{T/0.710}-1\right)=\left(2^{T/4.50}-1\right).
\end{equation*}

(d) Since $T\gg4.50$ is assumed, the $1$ may be neglected in the above equation, giving
\begin{equation*}
0.0120\cdot2^{T/0.710}=2^{T/4.50} ,
\end{equation*}
that is,
\begin{equation*}
0.0120=2^{T(1/4.50-1/0.710)} .
\end{equation*}
\begin{equation*}
T=\frac{\log0.0120}{\log2\cdot\left(\dfrac{1}{4.50}-\dfrac{1}{0.71}\right)}=5.38 ,
\end{equation*}
\begin{equation*}
T=5.38\times10^{9}\ \text{years} .
\end{equation*}

(e) $T$ is not in fact $\gg4.5\times10^{9}$ years, but it is $>0.710\times10^{9}$ years. Substituting the approximate value of $T$ (call it $T^*=5.38\times10^{9}$ years) into the term $2^{T/4.50}$ of the equation of part (c), in which the $1$ has not been neglected, one may calculate the $T$ occurring in the term $2^{T/0.71}$, and so obtain a better approximation to $T$; repeating the above calculation then gives an even better approximation.
\begin{equation*}
0.0120\left(2^{T/0.71}-1\right)=2^{T^*/4.50}-1 ,
\end{equation*}
\begin{equation*}
2^{T/0.71}-1=\frac{2^{5.38/4.50}-1}{0.0120}=107.5 .
\end{equation*}
Hence
\begin{equation*}
T=0.71\,\frac{\log92.2}{\log2}=4.80 .
\end{equation*}
Taking next $T^*=4.80$,
\begin{equation*}
2^{T/0.71}-1=\frac{2^{4.80/4.50}-1}{0.0120}=91.2 ,
\end{equation*}
\begin{equation*}
T=0.71\,\frac{\log92.2}{\log2}=4.63 .
\end{equation*}
One more iteration gives $T=4.58$. Hence the more accurate answer for $T$ lies between $4.5\times10^{9}$ and $4.6\times10^{9}$ years (both values are counted as correct).

\textbf{Part D: A spherical charge distribution}

(a) The charge density is
\begin{equation*}
\rho=\frac{Q}{(4/3)\pi R^3}=\frac{3Q}{4\pi R^3} .
\end{equation*}
By Gauss's law,
\begin{equation*}
E=\frac{Qr}{4\pi\varepsilon_0R^3}\qquad(r\le R),
\end{equation*}
\begin{equation*}
E=\frac{Q}{4\pi\varepsilon_0r^2}\qquad(r>R).
\end{equation*}

(b) The energy density is $(1/2)\varepsilon_0E^2$. For $r\le R$, the energy within a spherical shell of radius $r$ and thickness $\dd r$ is
\begin{equation*}
\dd E=\frac{1}{2}\varepsilon_0E^2\,4\pi r^2\,\dd r
=\frac{1}{2}\frac{Q^2}{4\pi\varepsilon_0R^6}r^4\,\dd r .
\end{equation*}
The energy stored inside the sphere is
\begin{equation*}
E_{\text{in}}=\frac{1}{2}\frac{Q^2}{4\pi\varepsilon_0R^6}\int_0^R r^4\,\dd r
=\frac{1}{40}\frac{Q^2}{\pi\varepsilon_0R} .
\end{equation*}
For $r>R$, the energy within a shell is
\begin{equation*}
\dd E=\frac{1}{2}\frac{Q^2}{4\pi\varepsilon_0r^4}r^2\,\dd r .
\end{equation*}
The energy in the region $r>R$ is
\begin{equation*}
E_{\text{out}}=\frac{1}{2}\frac{Q^2}{4\pi\varepsilon_0}
\int_R^\infty\frac{1}{r^2}\,\dd r
=\frac{1}{8}\frac{Q^2}{4\pi\varepsilon_0R} .
\end{equation*}
Hence the total energy associated with the charge is
\begin{equation*}
E=E_{\text{in}}+E_{\text{out}}
=\frac{3}{20}\frac{Q^2}{4\pi\varepsilon_0R} .
\end{equation*}
The result may also be obtained by summing the work done in assembling the charge; this is omitted here.

\textbf{Part E: The rotating copper ring}

The kinetic energy of rotation of the ring gradually decreases because of the heat losses due to the induced current. The horizontal component of the Earth's magnetic field that produces the induced current in the ring is
\begin{equation*}
B=44.5\times10^{-6}\cos64\degree\ \unitv{T} .
\end{equation*}
When the plane of the ring makes an angle $\theta=\omega t$ with the horizontal component of the Earth's field, the magnetic flux through the ring is
\begin{equation*}
\Phi=B\pi a^2\sin\omega t ,
\end{equation*}
where $a$ is the radius of the ring.

The instantaneous induced emf is
\begin{equation*}
\frac{\dd\Phi}{\dd t}=B\pi a^2\omega\cos\omega t .
\end{equation*}
The average thermal power dissipated in the resistance $R$ of the ring during one revolution is
\begin{equation*}
P=\frac{B^2\pi^2a^4\omega^2}{2R} .
\end{equation*}
The moment of inertia of the ring is $I=(1/2)ma^2$, where $m$ is the mass of the ring. The kinetic energy of rotation is therefore
\begin{equation*}
W=\frac{1}{2}I\omega^2=\frac{1}{4}ma^2\omega^2 .
\end{equation*}
Hence the energy equation is
\begin{equation*}
\frac{\dd W}{\dd t}=-P ,
\end{equation*}
that is,
\begin{equation*}
\frac{1}{2}ma^2\omega\frac{\dd\omega}{\dd t}
=-\frac{B^2\pi^2a^4\omega^2}{2R} ,
\end{equation*}
or
\begin{equation*}
\frac{\dd\omega}{\omega}=-\frac{B^2\pi^2a^2}{mR}\,\dd t .
\end{equation*}
If $T$ is the time required for the angular velocity to fall to half its value, then
\begin{equation*}
\int_\omega^{\omega/2}\frac{\dd\omega}{\omega}
=-\int_0^T\frac{B^2\pi^2a^2}{mR}\,\dd t ,
\end{equation*}
that is,
\begin{equation*}
\ln2=\frac{B^2\pi^2a^2}{mR}T .
\end{equation*}
But $R=2\pi a\rho/A$ and $m=2\pi aAd$, where $A$ is the cross-sectional area of the copper wire, $d$ is the density of copper and $\rho$ is its resistivity. Substituting these into the above gives
\begin{equation*}
T=\frac{4\rho d\ln2}{B^2} .
\end{equation*}
Inserting the data,
\begin{equation*}
T=1.10\times10^6\unitv{s}=306\unitv{hr}=12\unitv{days}\ 18\unitv{hr} .
\end{equation*}
This part may also be solved by a dynamical method; this is omitted here.

\subsection*{Problem 2}

(a) A cathode-ray tube consists of an electron gun and a fluorescent screen and is placed in a uniform, constant magnetic field of flux density $B$, as shown in Fig.~31.1. The direction of the magnetic field is parallel to the axis of the electron beam of the gun. The electron beam emerges from the anode of the gun along the axis, but with a divergence angle of at most $5\degree$ with the axis, as shown in Fig.~31.2. Normally the electron beam produces a diffuse spot on the fluorescent screen, but for a certain magnitude of the magnetic field a well-focused spot is obtained. Consider an electron which leaves the gun with its direction of motion at an angle $\beta$ $(0<\beta<5\degree)$ to the axis, and resolve its motion into components perpendicular and parallel to the axis. Using the following quantities, derive an expression for the specific charge $e/m$ of the electron: the minimum magnetic flux density that focuses the beam to a point; the accelerating potential difference $V$ of the gun (note that $V<2\unitv{kV}$); the distance $D$ from the anode to the screen.

\begin{figbox}{31.1}
\begin{tikzpicture}[scale=0.9]
\draw[thickline] (-3.2,-0.8) -- (1.9,-0.8) -- (1.9,0.8) -- (-3.2,0.8) -- cycle;
\draw[thickline] (-2.7,-0.32) rectangle (-1.9,0.32);
\node at (-2.3,0) {gun};
\begin{scope}
\clip (1.9,-0.8) rectangle (2.15,0.8);
\foreach \y in {-0.8,-0.6,...,0.8} \draw[thin] (1.9,\y) -- (2.15,\y+0.3);
\end{scope}
\draw[thickline] (1.9,-0.8) rectangle (2.15,0.8);
\node[right] at (2.15,0) {screen};
\draw[->,thickline] (-1.6,-1.45) -- (0.6,-1.45) node[right] {$B$};
\end{tikzpicture}
\end{figbox}

\begin{figbox}{31.2}
\begin{tikzpicture}[scale=0.9]
\draw[thickline] (-3.3,-0.4) rectangle (-2.5,0.4);
\node at (-2.9,0) {gun};
\draw[dashedline] (-2.5,0) -- (2.0,0);
\draw[thickline] (-2.4,0.06) -- (1.8,0.72);
\draw[thickline] (-2.4,-0.06) -- (1.8,-0.72);
\draw[thickline] (1.9,-0.75) -- (1.9,0.75);
\node at (1.0,0.22) {$5\degree$};
\node at (1.0,-0.22) {$5\degree$};
\end{tikzpicture}
\end{figbox}

(b) Consider another method of measuring the specific charge of the electron. Fig.~31.3 shows the side view and the plan view of the experimental arrangement, and also the direction of the magnetic field $B$. Two copper discs separated by a very small distance $t$ are placed in this uniform magnetic field; the radius of the discs is $\rho$. A potential difference $V$ is maintained between the two discs. The discs are parallel and coaxial, and their axis is perpendicular to the magnetic field. A photographic plate is placed coaxially on a cylinder of radius $\rho+s$, on its inner surface; in other words, the plate is at a radial distance $s$ from the edge of the discs. The whole apparatus is placed in vacuum. Note that $t$ is much smaller than $s$ and $\rho$.

A point source of $\beta$ particles is placed between the centres of the two discs and emits $\beta$ particles uniformly in all directions, the speeds of the particles being distributed over a certain range. The same plate is exposed under the following different experimental conditions:

\begin{center}
\begin{tabular}{ll}
Case 1 & $B=0$, $V=0$;\\
Case 2 & $B=B_0$, $V=V_0$;\\
Case 3 & $B=-B_0$, $V=-V_0$.
\end{tabular}
\end{center}

Here $V_0$ and $B_0$ are positive constants. Note that when $V>0$ the upper disc is positively charged (when $V<0$ the upper disc is negatively charged). When $B>0$ the direction of the magnetic field is defined as shown in Fig.~31.3 (when $B<0$ the field is in the opposite direction). In solving this part the separation of the two discs may be neglected.

In Fig.~31.3 the two regions of the plate are marked $A$ and $B$ respectively. After the plate has been exposed and developed, the $\beta$-particle tracks displayed by the developed plate in one of the two regions are as shown in Fig.~31.4. From which region ($A$ or $B$) was this part of the plate taken? Justify your answer by indicating the direction of the force acting on the electron.

\begin{figbox}{31.3}
\begin{tikzpicture}[scale=0.85]
% side view
\draw[thickline] (-1.0,0.35) -- (-1.0,1.35);
\draw[thickline] (3.9,0.35) -- (3.9,1.35);
\node[above] at (-1.0,1.35) {plate};
\node[above] at (3.9,1.35) {plate};
\draw[thickline] (0,0.62) -- (3.4,0.62);
\draw[thickline] (0,0.86) -- (3.4,0.86);
\node[above] at (1.5,0.86) {disc};
\draw[<->,thin] (-0.55,0.62) -- (-0.55,0.86);
\node[left] at (-0.6,0.74) {$t$};
\draw[<->,thin] (1.7,0.42) -- (3.4,0.42);
\node[below] at (2.55,0.42) {$\rho$};
\draw[<->,thin] (3.4,0.42) -- (3.9,0.42);
\node[below] at (3.7,0.38) {$s$};
\draw[->,thickline] (0.6,0.15) -- (1.8,0.15) node[right] {$B$};
% plan view
\begin{scope}[yshift=-3.6cm]
\draw[thickline] (0,0) circle (1.0);
\draw[thickline] (0,0) circle (1.55);
\fill (0,0) circle (1.5pt);
\node at (-0.35,0.45) {disc};
\draw[->] (0,0) -- (35:1.0);
\node at (18:0.55) {$\varphi$};
\draw[->] (0,0) -- (0:1.55);
\node[right] at (1.55,0.15) {plate};
\draw[<->,thin] (0,0.0) -- (0,1.0);
\node[left] at (-0.08,0.55) {$\rho$};
\draw[<->,thin] (1.0,0.0) -- (1.55,0.0);
\node[below] at (1.28,-0.05) {$s$};
\node at (2.1,0.75) {region $A$};
\node at (2.1,-0.75) {region $B$};
\draw[->,thickline] (-1.2,-1.75) -- (0.0,-1.75) node[right] {$B$};
\node at (0,-2.1) {plan view};
\end{scope}
\end{tikzpicture}
\end{figbox}

\begin{figbox}{31.4}
\begin{tikzpicture}[scale=0.9]
\draw[thickline] (-3.0,-1.2) rectangle (3.8,1.2);
\foreach \k in {-0.9,-0.6,-0.3,0,0.3,0.6,0.9}{
  \draw[thin] plot[domain=-3.0:3.8,samples=80] (\x,{\k+0.12*sin(60*\x)});
}
\draw[<->,thin] (-2.7,-0.9) -- (-2.7,0.9);
\node[left] at (-2.8,0) {$y$};
\node[below] at (-2.2,-1.25) {$\varphi=0\degree$};
\node[below] at (0.2,-1.25) {$\varphi=90\degree$};
\node[below] at (2.9,-1.25) {$\varphi=180\degree$};
\end{tikzpicture}
\end{figbox}

(c) After exposure and development, the plate, unrolled, is as shown in Fig.~31.4. The separation of the two outermost tracks is measured with a microscope [the separation $y$ at a particular angle is marked in Fig.~31.4]. The results of these measurements are given in the table below. The angle $\varphi$ is defined as the angle between the direction of the magnetic field and the line joining a point on the plate to the centre of the discs, as shown in Fig.~31.3.

\begin{center}
Table 31.1
\end{center}
\begin{center}
\begin{tabular}{|c|c|c|c|c|c|c|c|}
\hline
Angle with magnetic field/$^\circ$ & $\varphi$ & 90 & 60 & 50 & 40 & 30 & 23\\
\hline
Separation/mm & $y$ & 17.4 & 12.7 & 9.7 & 6.4 & 3.3 & end of track\\
\hline
\end{tabular}
\end{center}

The numerical values of the parameters of the measuring system are as follows:
\begin{equation*}
B_0=6.91\unitv{mT}\qquad V_0=580\unitv{V}\qquad
t=0.80\unitv{mm}\qquad s=41.0\unitv{mm} .
\end{equation*}
In addition, the speed of light in vacuum may be taken as $3.00\times10^8\unitv{m\cdot s^{-1}}$ and the rest mass of the electron as $9.11\times10^{-31}\unitv{kg}$. Calculate the maximum kinetic energy of the observed $\beta$ particles (in eV).

(d) Using the information given in part (c) above, find the ratio of the charge of the electron to its rest mass. This should be done by plotting a suitable graph on the graph paper provided. Write down the algebraic expressions for the quantities represented by the horizontal and vertical axes of your graph, and give the value of the specific charge $e/m$.

Note that, because of systematic errors of observation, the value you obtain may differ from the accepted standard value.

\subsection*{Problem 2}

(a) According to the problem, when an electron completes exactly one revolution just as it reaches the screen, it will be sharply focused on the screen.

The angular velocity of the circular motion is $\omega=eB/m$, and the period of one revolution is $T=2\pi m/eB$.

The speed of the electron is $u=\sqrt{2eV/m}$ (since $V<2\unitv{kV}$, it is not difficult to calculate that $u<2.65\times10^7\unitv{m\cdot s^{-1}}$, so relativistic effects need not be considered).

The distance travelled by the electron is
\begin{equation*}
D=Tu\cos\beta\approx Tu=(2\pi m/eB)\sqrt{2eV/m} .
\end{equation*}
Hence the specific charge of the electron is
\begin{equation*}
e/m=8\pi^2V/B^2D^2 .
\end{equation*}

(b) Consider Case 2. The electric force on the electron is directed upwards.

In region $A$ the magnetic force on the electron is also directed upwards, so the electron strikes the disc and cannot reach the plate.

In region $B$ the magnetic force is directed downwards; if the magnetic force equals the electric force, the electron experiences no net force and emerges from between the discs, exposing the plate.

Case 3 is similar. Hence the plate was taken from region $B$.

(c) An electron that emerges from between the discs is in force balance. The electric force on it is $eV/t$ and the magnetic force is $euB\sin\varphi$; hence
\begin{equation*}
eV/t=euB\sin\varphi ,
\end{equation*}
giving the electron speed
\begin{equation*}
u=V/Bt\sin\varphi .
\end{equation*}
The maximum speed $u_m$ corresponds to the smallest angle $\varphi$, namely $23\degree$; hence
\begin{align*}
u_m&=\frac{V_0}{B_0t\sin23\degree}
=\frac{580}{6.91\times10^{-3}\times0.80\times10^{-3}\times\sin23\degree}
\unitv{m\cdot s^{-1}}\\
&=2.685\times10^8\unitv{m\cdot s^{-1}}=0.895c .
\end{align*}
This is relativistic, so the maximum kinetic energy of the electron is
\begin{align*}
E_m&=(\gamma-1)mc^2
=\left(\frac{1}{\sqrt{1-u^2/c^2}}-1\right)mc^2\\
&=1.02\times10^{-13}\unitv{J}=637\unitv{keV} .
\end{align*}

(d) An electron emerging from between the discs is acted on only by the magnetic force, and this force is approximately vertical, because the angle between the velocity of the electron and the horizontal is always very small.

The acceleration due to the magnetic force is $a=Beu\sin\varphi/\gamma m$. The initial horizontal speed of the electron is $u$, and the time taken by the electron to travel from between the discs to the plate is $s/u$. During this time the vertical displacement is $y/2=(1/2)a(s/u)^2$, that is,
\begin{equation*}
y=\frac{Bes^2\sin\varphi}{\gamma mu} .
\end{equation*}
From (c), an electron emerging from between the discs must have $u=\dfrac{V}{Bt\sin\varphi}$. Substituting this to eliminate $u$ [note that $\gamma=1/\sqrt{1-u^2/c^2}$] gives
\begin{equation*}
y^2=\left(\frac{eBs\sin\varphi}{m}\right)^2
\left[\left(\frac{Bt\sin\varphi}{V}\right)^2-\left(\frac{s}{c}\right)^2\right].
\end{equation*}
Plotting $(y/Bs\sin\varphi)^2$ as ordinate against $(Bt\sin\varphi/V)^2$ as abscissa, the resulting straight line has slope $(e/m)^2$ and intercept on the ordinate axis $-(es/mc)^2$.

From the graph drawn (figure omitted) the intercept is read as $-537.7\unitv{(C\cdot s/kg)^2}$, giving $e/m=1.70\times10^{11}\unitv{C/kg}$.

The slope is read as $2.826\times10^{22}\unitv{(C/kg)^2}$, giving $e/m=1.68\times10^{11}\unitv{C/kg}$.

\subsection*{Problem 3}

\textbf{Gravitational waves and the effects of gravity on light}

\textbf{Part A}

This part is concerned with the difficulty of detecting the gravitational waves produced by astronomical events. It is known that a distant supernova explosion may produce a disturbance of about $10^{-19}\unitv{N\cdot kg^{-1}}$ in the gravitational field strength at the Earth's surface. A model of a gravitational-wave detector (see Fig.~31.5) consists of two metal bars each $1\unitv{m}$ long, mounted at right angles. One end of each bar is polished to an optical flat, and the other end is rigidly fixed. The position of one of the bars is adjusted so that the signal received by the phototube is a minimum (see Fig.~31.5).

A very short longitudinal pulse is produced in the bar by a piezoelectric device, with the result that the free end of the bar vibrates with a longitudinal displacement $\Delta x_1$,
\begin{equation*}
\Delta x_1=ae^{-\mu t}\cos(\omega t+\varphi),
\end{equation*}
where $a$, $\mu$, $\omega$ and $\varphi$ are constants.

\begin{enumerate}
\item[(a)] If the amplitude of the displacement decreases by $20\%$ in a time interval of $50$ seconds, find the value of $\mu$.
\item[(b)] Suppose that both bars are made of aluminium, with density $\rho=2700\unitv{kg\cdot m^{-3}}$ and Young's modulus $E=7.1\times10^{10}\unitv{Pa}$. Given that the speed of longitudinal waves is $v=\sqrt{E/\rho}$, find the minimum value of $\omega$.
\item[(c)] It is in general impossible to make the two bars exactly the same length, so that the phototube signal shows a beat frequency of $0.005\unitv{Hz}$. What is the difference in length of the two bars?
\item[(d)] For a bar of length $l$, derive an algebraic expression for the change in length $\Delta l$ caused by a change $\Delta g$ in the gravitational field strength $g$, in terms of $l$ and other constants of the bar material. Assume that the response of the detector to a change in gravitational field strength occurs only along the axis of one bar.
\item[(e)] A laser produces monochromatic light of wavelength $656\unitv{nm}$. If the smallest fringe shift that can be measured is $10^{-4}$ of the laser wavelength, what is the minimum value of $l$ for the bar if the system is to be able to detect a change in $g$ of $10^{-19}\unitv{N\cdot kg^{-1}}$?
\end{enumerate}

\begin{figbox}{31.5}
\begin{tikzpicture}[scale=0.95]
\draw[thickline] (-3.3,0.06) -- (0,0.06);
\draw[thickline] (-3.3,-0.06) -- (0,-0.06);
\fill (-3.75,-0.35) rectangle (-3.3,0.35);
\draw[thickline] (0.06,0) -- (0.06,-2.4);
\draw[thickline] (-0.06,0) -- (-0.06,-2.4);
\fill (-0.35,-2.75) rectangle (0.35,-2.4);
\draw[thickline] (-0.1,-3.05) rectangle (0.1,-2.75);
\node[below] at (0,-3.05) {adjusting screw};
\draw[thickline] (-0.42,0.42) -- (0.42,-0.42);
\node[above] at (0.55,0.55) {half-silvered mirror};
\draw[thickline] (1.8,-0.22) rectangle (3.3,0.22);
\node at (2.55,0) {laser};
\draw[thin] (1.8,0) -- (0.45,0);
\draw[thickline] (1.35,-1.25) rectangle (1.75,-0.85);
\node[right] at (1.75,-1.05) {phototube};
\draw[thin] (0.35,-0.35) -- (1.5,-0.97);
\draw[->,thin] (-2.9,0.45) -- (-2.4,0.45) node[above] {$x$};
\end{tikzpicture}
\end{figbox}

\textbf{Part B}

This part examines the effects produced by the gravitational field on the propagation of light in space.

(a) A photon emitted from the surface of the Sun (mass $M$, radius $R$) will be red-shifted. Assuming that the mass of a photon is equivalent to its energy, use Newton's theory of gravitation to show that the effective (or measured) frequency of the photon at infinity is reduced by a factor $(1-GM/Rc^2)$ (that is, it is red-shifted).

(b) A decrease in the frequency of light is equivalent to an increase in its period. When light is used as a standard clock, this is equivalent to a dilation of time. Moreover, a dilation of time is always accompanied by a contraction of length by the same factor.

We now attempt to study the effect of this on light propagating at the edge of the Sun. First define the equivalent refractive index $n_r$ at a distance $r$ from the centre of the Sun:
\begin{equation*}
n_r=c/c'_r ,
\end{equation*}
where $c$ is the speed of light measured in a coordinate system far from the gravitational influence of the Sun ($r\to\infty$), and $c'_r$ is the speed of light measured in the coordinate system at a distance $r$ from the centre of the Sun.

When $GM/rc^2$ is very small, show that $n_r$ may be expressed approximately as
\begin{equation*}
n_r=1+aGM/rc^2 ,
\end{equation*}
where $a$ is a constant; determine this constant.

(c) Using the above expression for $n_r$, calculate the angle (in radians) by which light is deflected from a straight-line path when passing the edge of the Sun.

Data:
\begin{align*}
\text{gravitational constant } G&=6.67\times10^{-11}\unitv{N\cdot m^2\cdot kg^{-2}},\\
\text{mass of the Sun } M&=1.99\times10^{30}\unitv{kg},\\
\text{radius of the Sun } R&=6.95\times10^{8}\unitv{m},\\
\text{speed of light } c&=3.00\times10^{8}\unitv{m\cdot s^{-1}} .
\end{align*}
You may need the following integral:
\begin{equation*}
\int_{-\infty}^{+\infty}\frac{\dd x}{(x^2+a^2)^{3/2}}=\frac{2}{a} .
\end{equation*}

\subsection*{Problem 3}

\textbf{Part A}

(a) Since
\begin{equation*}
\Delta x_1=ae^{-\mu t}\cos(\omega t+\varphi),
\end{equation*}
\begin{equation*}
0.8=\ee^{-50\mu}\quad\Longrightarrow\quad
\mu=-\frac{\ln0.8}{50}\unitv{s^{-1}}=4.5\times10^{-3}\unitv{s^{-1}} .
\end{equation*}

(b)
\begin{equation*}
v=(E/\rho)^{1/2}=(7.1\times10^{10}/2700)^{1/2}\unitv{m\cdot s^{-1}}
=5.1\times10^3\unitv{m\cdot s^{-1}} .
\end{equation*}
For the fundamental mode of vibration,
\begin{equation*}
\lambda=4l=4\unitv{m},
\end{equation*}
\begin{equation*}
v=\lambda f,\qquad f=v/\lambda ,
\end{equation*}
\begin{equation*}
\omega=2\pi f=2\pi v/\lambda=8.06\times10^3\unitv{rad\cdot s^{-1}} .
\end{equation*}

(c)
\begin{equation*}
f=v/\lambda=v/4l,\qquad |\Delta f|=v\Delta l/4l^2 ,
\end{equation*}
\begin{equation*}
\Delta l=4l^2\Delta f/v=4\times1^2\times0.05/5.1\times10^3\unitv{m}
=3.9\times10^{-6}\unitv{m} .
\end{equation*}

(d) At a distance $x$ from the free end, the additional gravitational force on the bar due to the change in gravitational field strength is
\begin{equation*}
m\Delta g=\rho xA\Delta g ,
\end{equation*}
where $A$ is the cross-sectional area of the bar. The corresponding additional stress is $m\Delta g/A=\rho x\Delta g$, and the additional strain produced by it is $\delta(\dd x)/\dd x=\rho x\Delta g/E$. Hence the additional deformation of the bar is
\begin{equation*}
\Delta l=\sum\delta(\dd x)=\int_0^l\frac{\rho\Delta g}{E}x\,\dd x
=\frac{\rho\Delta g}{2E}l^2 .
\end{equation*}

(e) When one bar is deformed by $\Delta l$ because of the change in gravitational field strength, the optical path difference of the laser light reflected from the free ends of the two bars changes by $2\Delta l$. According to the problem, $2\Delta l=10^{-4}\lambda$; substituting the expression for $\Delta l$ from (d) gives
\begin{equation*}
10^{-4}\lambda=(\rho\Delta g/E)l^2 .
\end{equation*}
Hence
\begin{align*}
l&=\sqrt{\frac{10^{-4}\lambda E}{\rho\Delta g}}
=\sqrt{\frac{10^{-4}\times656\times10^{-9}\times7.1\times10^{10}}
{2700\times10^{-19}}}\unitv{m}\\
&=1.31\times10^{8}\unitv{m} .
\end{align*}

\textbf{Part B}

(a) Since $mc^2=hf$, we have $m=hf/c^2$. Let the frequency of the photon at the surface of the Sun be $f_R$ and its frequency at infinity be $f$. By conservation of energy,
\begin{equation*}
hf=hf_R-GMm/R=hf_R-\frac{GMhf_R}{Rc^2}
=hf_R\left(1-\frac{GM}{Rc^2}\right).
\end{equation*}
Hence
\begin{equation*}
f=f_R\left(1-GM/Rc^2\right).
\end{equation*}

(b) Let the frequency of the photon at a distance $r$ from the centre of the Sun be $f_r$; in the same way,
\begin{equation*}
f=f_r\left(1-GM/rc^2\right).
\end{equation*}
The above relation shows that the frequency of a photon at infinity is lower than its frequency at a distance $r$ from the centre of the Sun by a factor, and its period is longer by the same factor. In other words, the frequency of a photon at a distance $r$ from the centre of the Sun is higher than the frequency of light at infinity by a factor, and its period is shorter by the same factor. When a photon is used as a standard clock, the time reading of the same process at a distance $r$ from the centre of the Sun will be greater than the time reading at infinity by a factor, and the length reading will be smaller by the same factor. Therefore the speed of light measured at a distance $r$ from the centre of the Sun is smaller than the speed of light measured at infinity by the square of this factor, that is,
\begin{equation*}
c'_r=c\left(1-GM/rc^2\right)^2 .
\end{equation*}
Hence, when $GM/rc^2\ll1$,
\begin{equation*}
n_r=c/c'_r=\left(1-GM/rc^2\right)^{-2}\approx1+2GM/rc^2 .
\end{equation*}
Therefore
\begin{equation*}
a=2 .
\end{equation*}

(c) Since the refractive index decreases with increasing distance $r$ from the centre of the Sun, light passing the edge of the Sun is bent towards the centre of the Sun by refraction (Fig.~31.6). By Snell's law,
\begin{equation*}
n(r+\delta r)\sin(\theta+\delta\theta)=n(r)\sin\theta ,
\end{equation*}
that is,
\begin{equation*}
\left[n(r)+\frac{\dd n}{\dd r}\delta r\right]
\left(\sin\theta+\cos\theta\cdot\delta\theta\right)=n(r)\sin\theta ,
\end{equation*}
\begin{equation*}
\sin\theta\,\frac{\dd n}{\dd r}\delta r+n(r)\cos\theta\cdot\delta\theta=0 ,
\end{equation*}
\begin{equation*}
\delta\theta=-\tan\theta\cdot\frac{1}{n}\cdot\frac{\dd n}{\dd r}\delta r .
\end{equation*}
But
\begin{equation*}
n=1+2GM/rc^2\approx1,\qquad
\frac{\dd n}{\dd r}=-2GM/r^2c^2 .
\end{equation*}
Substituting into the above,
\begin{equation*}
\delta\theta=\tan\theta\left(2GM/r^2c^2\right)\delta r .
\end{equation*}
From the figure, $r^2=x^2+R^2$, $r\,\dd r=x\,\dd x$ and $\tan\theta=R/x$; hence the deflection angle is
\begin{align*}
\Delta\theta&=\int\delta\theta
=\frac{2GM}{c^2}\int\frac{\tan\theta\,\dd r}{r^2}
=\frac{2GM}{c^2}\int\frac{R}{x}\frac{x\,\dd x}{r^3}\\
&=\frac{2GMR}{c^2}\int_{-\infty}^{+\infty}
\frac{\dd x}{(x^2+R^2)^{3/2}}
=\frac{2GMR}{c^2}\cdot\frac{2}{R^2}
=\frac{4GM}{Rc^2} .
\end{align*}
Inserting the data,
\begin{equation*}
\Delta\theta=\frac{4\times6.67\times10^{-11}\times1.99\times10^{30}}
{6.95\times10^8\times(3\times10^8)^2}\unitv{rad}
=8.5\times10^{-6}\unitv{rad} .
\end{equation*}

\begin{figbox}{31.6}
\begin{tikzpicture}[scale=0.95]
\draw[thickline] (1.5,-1.1) circle (1.1);
\node[below right] at (2.4,-1.8) {Sun};
\fill (1.5,-1.1) circle (1.2pt) node[below] {$O$};
\draw[thickline] plot[domain=-4.2:3.0,samples=60] (\x,{0.35+0.03*\x+0.02*\x*\x});
\draw[thickline] (1.5,-1.1) -- (-0.8,0.05);
\node[below] at (0.35,-0.55) {$r$};
\draw[thickline] (1.5,-1.1) -- (1.5,0.12);
\node[right] at (1.53,-0.5) {$R$};
\draw[dashedline] (-0.8,0.05) -- (1.5,0.05);
\node[above] at (0.3,0.06) {$x$};
\draw[thickline] (1.5,-1.1) -- (-1.6,0.29);
\node[above] at (-0.2,-0.3) {$r+\delta r$};
\draw[thin] (-0.8,0.05) ++(0:0.55) arc (0:28:0.55);
\node at (-0.15,0.22) {$\theta$};
\draw[thin] (-1.6,0.29) ++(0:0.4) arc (0:27:0.4);
\node at (-1.05,0.43) {$\theta+\delta\theta$};
\node[above] at (2.8,0.6) {light ray};
\end{tikzpicture}
\end{figbox}

\subsection*{Experimental Problems}

\subsection*{Problem 4}

\textbf{The CDROM spectrometer}

The aim of this experiment is to represent graphically how the conductance of a photoresistor varies with wavelength over the visible spectrum. You are not required to state the uncertainty of your measurements.

$^*$ Conductance $G=1/\text{resistance}$ (unit: siemens, $1\unitv{S}=1\unitv{\Omega^{-1}}$).

\textbf{[Contents of the experiment]}

The experiment consists of the following five parts:

\begin{enumerate}
\item[(1)] Use a concave reflection grating (made from a piece of CDROM) to produce a focused first-order spectrum from lamp $A$ (a $12\unitv{V}$, $50\unitv{W}$ tungsten filament lamp).
\item[(2)] Measure and plot the conductance of the photoresistor as a function of wavelength within this first-order spectrum.
\item[(3)] Show that the light emission of the filament of lamp $A$ behaves approximately as an ideal black body.
\item[(4)] Find the temperature of the lamp filament when it is connected to a $12\unitv{V}$ supply.
\item[(5)] After calculating the energy distribution of the spectrum emitted by lamp $A$, correct the curve of conductance against wavelength.
\end{enumerate}

\textbf{[Precautions]}

\begin{enumerate}
\item[(1)] Beware of burns.
\item[(2)] Lamp $B$ must not be connected to any potential difference exceeding $2.0\unitv{V}$.
\item[(3)] The resistance range of the multimeter must not be used on a circuit that is switched on.
\end{enumerate}

\textbf{[Procedure]}

(a) The apparatus shown in Fig.~31.7 is already set up, so that the light from lamp $A$ is incident normally on the grating and the photoresistor is placed in the focal plane of the first-order spectrum. Move the photoresistor along the first-order spectrum and observe how its resistance (measured with the multimeter on the $\times$ range) varies with position.

\begin{figbox}{31.7}
\begin{tikzpicture}[scale=0.8]
% panel (a)
\draw[thickline] (-4.3,0.5) circle (0.75);
\fill (-3.9,0.5) circle (1.5pt) node[above left] {$A$};
\node[below] at (-4.3,-0.35) {lamp $A$};
\draw[thin] (-3.7,0.75) -- (1.5,1.15);
\draw[thin] (-3.7,0.25) -- (1.5,0.15);
\draw[thickline] (1.55,0.1) .. controls (1.15,0.65) .. (1.55,1.2);
\node[right] at (1.6,0.65) {grating};
\node[right] at (1.6,0.25) {(CDROM strip)};
\draw[thin] (1.5,1.15) -- (-1.4,-1.3);
\draw[thin] (1.5,0.15) -- (-1.4,-1.3);
\draw[thickline,rotate around={-35:(-1.4,-1.3)}] (-1.55,-1.42) rectangle (-1.25,-1.18);
\node[left] at (-1.6,-1.3) {photoresistor};
\draw[thin] (-1.55,-1.45) -- (-2.4,-1.9);
\draw[thin] (-1.25,-1.45) -- (-2.1,-2.0);
\draw[thickline] (-3.1,-2.2) rectangle (-2.1,-1.7);
\node at (-2.6,-1.95) {$\Omega$};
\node[below] at (-2.6,-2.25) {ohmmeter};
\draw[thin] (0.15,0.62) arc (150:200:0.55);
\node at (0.45,0.95) {$\theta$};
\draw[<->,thin] (-3.7,1.55) -- (1.5,1.55);
\node[above] at (-1.1,1.55) {$10\unitv{cm}$};
\draw[thickline] (-4.6,-2.6) -- (2.2,-2.6);
\node[below] at (1.4,-2.6) {table};
\node at (-1.0,-3.0) {(a) part of the experimental arrangement};
% panel (b)
\begin{scope}[xshift=5.6cm,yshift=-0.2cm]
\draw[thickline] (-1.3,0.9) rectangle (1.3,1.25);
\fill[pattern=north east lines] (-0.6,0.95) rectangle (0.6,1.2);
\node[below] at (-0.9,0.9) {CDROM strip};
\node[below] at (0.75,0.9) {plastic strip};
\draw[thickline] (-1.0,-0.2) rectangle (1.0,0.6);
\fill[pattern=north east lines] (0.1,0.0) rectangle (0.95,0.55);
\node[right] at (1.0,0.28) {photoresistor};
\draw[thin] (-1.0,0.0) -- (-1.7,-0.3);
\draw[thin] (-1.0,0.3) -- (-1.7,0.0);
\draw[thickline] (-2.6,-0.55) rectangle (-1.7,-0.05);
\node at (-2.15,-0.3) {$\Omega$};
\node[below] at (0,-0.35) {(b) photoresistor and multimeter};
\end{scope}
\end{tikzpicture}
\end{figbox}

(b) (1) Measure and record the resistance $R$ of the photoresistor at different positions in the first-order spectrum. Record the data in the blank table provided.

(2) Using the graph paper provided, plot the conductance $G$ of the photoresistor against the wavelength $\lambda$.

[Note] The angle $\theta$ between the direction of light of wavelength $\lambda$ in the first-order spectrum and the direction in which white light is reflected from the grating is given by
\begin{equation*}
\sin\theta=\frac{\lambda}{d},
\end{equation*}
where $d$ is the grating spacing. This grating has 620 lines per millimetre.

Since the emission characteristics of the lamp have not been taken into account, the curve drawn in (b)(2) cannot correctly represent the relation between the conductance of the photoresistor and wavelength. These characteristics will be studied in parts (c) and (d), so that in part (e) a corrected curve is obtained.

[Note] In part (c) three multimeters are used as ammeters. Do not adjust or move these meters. Use the fourth multimeter (marked $\times$) to measure all voltages.

(c) If the filament of the $50\unitv{W}$ lamp radiates as a black body, it can be shown that the relation between the voltage $V$ applied across it and the current $I$ through it may be written as
\begin{equation*}
V^3=CI^5 ,
\end{equation*}
where $C$ is a constant. Measure the corresponding values of $V$ and $I$ for lamp $A$ (in its housing). The ammeter has already been connected; do not adjust it.

\begin{enumerate}
\item[(1)] Record the data and calculated values in the table on the answer sheet.
\item[(2)] Draw a suitable graph on the graph paper provided to show that the filament radiates as a black body.
\end{enumerate}

(d) In order to correct the curve in (b)(2) we need to know the working temperature of the tungsten filament of lamp $A$. This can be found from the variation of the resistance of the filament with temperature.

[Curve provided] You are provided with a curve of the resistivity of tungsten (in $\unitv{\mu\Omega\cdot cm}$) against temperature (in K) (Fig.~31.8).

\begin{figbox}{31.8}
\begin{tikzpicture}[scale=0.78]
% panel (a): resistivity versus temperature
\draw[axis] (0,0) -- (5.2,0) node[below] {temperature/K};
\draw[axis] (0,0) -- (0,4.2) node[above,rotate=90,yshift=6pt] {resistivity/$\mu\Omega\cdot$cm};
\foreach \x/\t in {0/0,1.0/500,1.7/1000,2.4/1500,3.1/2000,3.8/2500,4.5/3000,5.0/3500}
  \draw (\x,0.05) -- (\x,-0.05) node[below] {\small\t};
\foreach \y/\t in {0/0,0.7/20,1.4/40,2.1/60,2.8/80,3.5/100,4.1/120}
  \draw (0.05,\y) -- (-0.05,\y) node[left] {\small\t};
\draw[thickline] (0.3,0.2) -- (4.9,4.0);
\fill (0.3,0.2) circle (1.5pt);
\node[above right] at (0.32,0.22) {\small $300,\ 5.6$};
\node at (2.6,-0.9) {(a) resistivity of tungsten versus temperature};
% panel (b): circuit
\begin{scope}[xshift=6.4cm]
\draw[thickline] (0,0) rectangle (1.0,4.0);
\node[above] at (0.5,4.0) {AC supply};
\foreach \y/\t in {0.2/0 V,0.9/2 V,1.6/4 V,2.3/6 V,3.0/8 V,3.5/10 V,3.9/12 V}
  {\draw (1.0,\y) -- (1.35,\y); \node[right] at (0.05,\y) {\small\t};}
\draw[thickline] (1.35,2.3) -- (2.0,2.3) -- (2.0,3.1);
\draw[thickline] (2.0,3.1) -- (2.9,3.1);
\draw[thickline] (2.9,3.1) circle (0.3);
\draw[thin] (2.75,3.0) -- (3.05,3.2);
\node[right] at (3.2,3.1) {lamp $C$};
\draw[thickline] (3.2,3.1) -- (4.0,3.1) -- (4.0,0.2) -- (1.35,0.2);
\draw[thickline] (2.6,0.2) -- (3.4,0.2);
\draw[thickline] (2.95,0.2) circle (0.25);
\draw[thin] (2.85,0.13) -- (3.05,0.27);
\node[below] at (2.95,-0.05) {lamp $B$};
\node at (2.0,-0.9) {(b) meters not shown};
\end{scope}
\end{tikzpicture}
\end{figbox}

If the resistance of lamp $A$ at some known temperature can be found, its temperature when working from a $12\unitv{V}$ supply can be found from its resistance at that working voltage. Unfortunately its resistance at room temperature is too small to be measured accurately with this apparatus. However, you are provided with another, smaller lamp $C$, which has a larger resistance at room temperature that can be measured. According to the following procedure, lamp $C$ can be used as an intermediary. You are also provided with another $12\unitv{V}$, $50\unitv{W}$ lamp $B$ identical to lamp $A$. Lamps $B$ and $C$ are mounted on a board and connected as shown in Fig.~31.8.

\begin{enumerate}
\item[(1)] With lamp $C$ unlit, measure its resistance at room temperature (using multimeter $\times$, and taking room temperature to be $300\unitv{K}$). Record this resistance $R_{C1}$ on the answer sheet.
\item[(2)] Compare the filaments of lamps $B$ and $C$ using the circuit shown in Fig.~31.8. Vary the current through lamp $C$ with the variable resistor until you see the two superimposed filaments at the same temperature. If the small filament is at a lower temperature than the large one, it appears as a thin dark ring. When equal temperatures are reached, measure the resistances of lamps $B$ and $C$ and record their values $R_{C2}$ and $R_B$ on the answer sheet. Remember that the ammeter has already been connected.
\item[(3)] Use the resistance--temperature curve provided to find the temperature at which lamps $B$ and $C$ have the same brightness. Record this temperature $T_{2V}$ on the answer sheet.
\item[(4)] Measure the resistance of lamp $A$ (in its housing) when it is connected to the $12\unitv{V}$ AC supply. Again, the ammeter has already been connected; do not adjust it. Record this $R_{12V}$ in the table on the answer sheet.
\item[(5)] Using the resistances of lamp $A$ (in its housing) at $2\unitv{V}$ and at $12\unitv{V}$, and its temperature at $2\unitv{V}$, find the temperature at which it works when connected to the $12\unitv{V}$ supply. Record this temperature $T_{12V}$ in the table on the answer sheet.
\end{enumerate}

[Curve provided] You are provided with curves of the relative radiant intensity of a black body (Planck curves) at the temperatures $2000\unitv{K}$, $2250\unitv{K}$, $2500\unitv{K}$, $2750\unitv{K}$, $3000\unitv{K}$ and $3250\unitv{K}$ (Fig.~31.9).

(e) Using these graphs and the result of (d)(5), draw on the graph paper provided a corrected curve of the conductance of the photoresistor (in arbitrary units) against wavelength. Assume that the conductance of the photoresistor at any wavelength is proportional to the radiant intensity at that wavelength (this assumption is reasonable in this experiment, in which the radiant intensity falling on the photoresistor is low). Also assume that the grating diffracts uniformly over the first-order spectrum.

\begin{figbox}{31.9}
\begin{tikzpicture}[scale=0.72]
% panel (a)
\draw[axis] (0,0) -- (6.2,0) node[below] {wavelength/m};
\draw[axis] (0,0) -- (0,5.0);
\node[rotate=90,above] at (-0.55,2.5) {\small radiant intensity per unit wavelength};
\foreach \x/\t in {0/0.00E+00,1.5/2.00E-07,3.0/4.00E-07,4.5/6.00E-07,6.0/8.00E-07}
  \draw (\x,0.05) -- (\x,-0.05) node[below,font=\tiny] {\t};
\foreach \y/\t in {0/0.00E+00,1.2/2.00E+02,2.4/4.00E+02,3.6/6.00E+02,4.8/8.00E+02}
  \draw (0.05,\y) -- (-0.05,\y) node[left,font=\tiny] {\t};
\draw[thickline] (0,0) .. controls (3.0,0.15) and (4.6,1.5) .. (6.0,4.4);
\draw[thickline] (0,0) .. controls (3.2,0.1) and (4.9,0.9) .. (6.0,2.8);
\draw[thickline] (0,0) .. controls (3.4,0.05) and (5.1,0.5) .. (6.0,1.7);
\node at (3.7,3.9) {\small $2500\unitv{K}$};
\node at (4.5,2.3) {\small $2250\unitv{K}$};
\node at (5.0,1.25) {\small $2000\unitv{K}$};
\node at (3.0,-1.0) {\small (a) Planck black-body curves ($2000\unitv{K}$, $2250\unitv{K}$, $2500\unitv{K}$)};
% panel (b)
\begin{scope}[yshift=-8.0cm]
\draw[axis] (0,0) -- (6.2,0) node[below] {wavelength/m};
\draw[axis] (0,0) -- (0,5.0);
\node[rotate=90,above] at (-0.55,2.5) {\small radiant intensity per unit wavelength};
\foreach \x/\t in {0/0.00E+00,1.5/2.00E-07,3.0/4.00E-07,4.5/6.00E-07,6.0/8.00E-07}
  \draw (\x,0.05) -- (\x,-0.05) node[below,font=\tiny] {\t};
\foreach \y/\t in {0/0.00E+00,1.2/1.00E+03,2.4/2.00E+03,3.6/3.00E+03,4.8/4.00E+03}
  \draw (0.05,\y) -- (-0.05,\y) node[left,font=\tiny] {\t};
\draw[thickline] (0,0) .. controls (3.0,0.15) and (4.5,1.8) .. (5.8,4.7);
\draw[thickline] (0,0) .. controls (3.2,0.1) and (4.8,1.1) .. (6.0,3.3);
\draw[thickline] (0,0) .. controls (3.4,0.05) and (5.1,0.6) .. (6.0,2.1);
\node at (3.9,4.4) {\small $3250\unitv{K}$};
\node at (4.6,2.9) {\small $3000\unitv{K}$};
\node at (5.1,1.6) {\small $2750\unitv{K}$};
\node at (3.0,-1.0) {\small (b) Planck black-body curves ($2750\unitv{K}$, $3000\unitv{K}$, $3250\unitv{K}$)};
\end{scope}
\end{tikzpicture}
\end{figbox}

\subsection*{Solutions to the Experimental Problems}

\subsection*{Problem 4}

(a) Connect up the apparatus and instruments as required in the figure. Let the white light produced by lamp $A$ fall normally on the concave reflection grating, and observe the diffracted light of different wavelengths.

(b) Move the photoresistor from position to position, measure the relation between the wavelength $\lambda$ (in nm) of the light in the first-order spectrum diffracted by the grating and the resistance $R$ of the photoresistor, and convert the resistance $R$ into the conductance $G=1/R$, entering it in the table.

The wavelength of the light is found from the formula $\sin\theta=\lambda/d$; the angle $\theta$ can be obtained by measuring lengths on the base plate with the paper and ruler provided, and then calculating. The resistance of the photoresistor is measured with the ohmmeter range of the digital multimeter marked $\times$. The measured range of wavelengths is the visible light from 400 to 700 nm. There should be more than 9 data points; note that near the peak of the curve the data points should be closer together. Plot the curve of the conductance $G$ of the photoresistor against the wavelength $\lambda$ on the graph paper; the peak of this curve lies in the range $(580\pm20)\unitv{nm}$. (Figure omitted.)

(c) Connect the current range of the digital multimeter in series in the circuit of lamp $A$, and use the voltage range of the multimeter marked $\times$ to measure the voltage across lamp $A$. Measure the relation between the voltage $V$ across the lamp and the current $I$ when lamp $A$ is lit. Plotting $V^3$ against $I^5$ gives a straight line, which shows that the filament of lamp $A$ radiates as a black body.

(d) Measuring the filament resistance $R_C$ of lamp $C$, not connected to the supply, at room temperature $300\unitv{K}$ gives $R_{C1}=13.5\unitv{\Omega}$; recording next the resistance $R_{C2}$ of lamp $C$ when it is as bright as lamp $B$ gives $R_{C2}=75.0\unitv{\Omega}$ (measured by the volt--ampere method). Hence, while measuring the voltage across lamp $B$, the resistance $R_B=1.20\unitv{\Omega}$ of lamp $B$ when it emits the same intensity of light can be found. Consulting curve 1, namely the data for the resistivity of tungsten against temperature $T$, calculation gives the filament temperature of lamps $C$ and $B$ when they have the same brightness as $T_{2V}=2130\unitv{K}$.

Using the volt--ampere method, the resistance of lamp $A$ when connected to the $12\unitv{V}$ AC supply is found to be $R_{12V}=2.85\unitv{\Omega}$. Similarly, consulting the curve of the resistivity of tungsten against temperature $T$, from $R_B$, $T_{2V}$ and $R_{12V}$ one obtains the filament temperature of lamp $A$ when connected to the $12\unitv{V}$ AC supply, $T_{12V}=2940\unitv{K}$, close to $3000\unitv{K}$.

(e) According to the assumption, the conductance of the photoresistor at any wavelength is proportional to the radiant intensity at that wavelength. Consult the $3000\unitv{K}$ curve of Fig.~31.9(b), the Planck black-body radiation curve (radiant intensity per unit wavelength against wavelength). Correct the measured conductance $G$ of the photoresistor to obtain the corrected conductance $G'$, and plot the $G'$--$\lambda$ relation. (Figure omitted.)

\subsection*{Problem 5}

\textbf{The magnetic slider}

In this experiment, the uncertainties must be stated for all measured data and experimental results, and the data points in the graphs must also be marked with their uncertainties.

\textbf{[Aim]}

To study the forces acting on a slider as it slides down an inclined plane.

\textbf{[Precautions]}

Do not touch the round face of the slider or the paper surface on the inclined track directly with your hands. Use the gloves provided. For convenience, the two faces of the round slider have been covered with paper of different colours.

\textbf{[Timing]}

Two magnetic sensors are mounted below the inclined plane of the sliding track and are used to trigger the circuit of the ``time-measuring box'' on the table. When the slider moves between the two magnetic sensors, the green light-emitting diode on the box lights up. The multimeter is used to measure the potential difference across a capacitor; when the green light-emitting diode is lit, the capacitor is connected to a constant-current source and charged (its charging current is constant and proportional to the potential difference across the battery). Hence the reading of the charged voltage displayed by the multimeter can be used to measure the time taken by the slider to travel between the two magnetic sensors. This reading can already represent, in some unit, the speed of the slider between the two sensors.

\textbf{[Operating the timer]}

(1) While the black button on the side of the time-measuring box is pressed, the circuit of the box is switched on and can be used for time measurement.

(2) At this time, if the green light is on, place the slider with its light-coloured face upwards (in this experiment the light-coloured face of the slider is always kept upwards) and put it between the two sensors; let it slide down the incline past the lower sensor, whereupon the green light should go out.

(3) Before releasing the slider, the potential difference across the capacitor can be discharged to zero by pressing and holding the red button on the timer box (for at least 10 seconds).

(4) The battery voltage can be measured by connecting the multimeter to the two terminals marked with the battery symbol.

\textbf{[Definitions]}

(1) A body sliding down an inclined plane is acted on by a resistive force $F$ along the plane and a normal force $N$. Define
\begin{equation*}
\xi=\frac{F}{N} .
\end{equation*}

(2) When the resistance arises only from friction, $\xi$ is equal to $\mu_k$, which is called the coefficient of kinetic friction of the surface; it is independent of speed.

(3) When the blue (dark) face is in contact with the inclined plane, define
\begin{equation*}
\xi_d=\frac{F_d}{N} .
\end{equation*}
Here the force $F_d$ along the plane arises partly from surface friction and partly from magnetic effects.

(4) Define the variable corresponding to magnetic effects alone:
\begin{equation*}
\xi_{d'}=\xi_d-\mu_k .
\end{equation*}

\textbf{[Important hints and suggestions]}

(1) It is very useful at the beginning to study the sliding behaviour of the slider qualitatively.

(2) Before making quantitative measurements, think about the underlying physics. Remember to present the results by means of graphs as far as possible.

(3) Unless you have enough time, do not record too much experimental data.

(4) In the experiment you measure the potential difference across an electrolytic capacitor; unlike a simple air capacitor, a slow leakage is normal, and the potential difference will not remain perfectly stable.

(5) You are provided with a slider and a $9.0\unitv{V}$ battery; please economise on electricity. The constant current that charges the capacitor is proportional to the potential difference across the battery; note the potential difference across the battery and record the result. If the potential difference across the battery falls below $8.4\unitv{V}$, the sensors may not work reliably, and the battery should then be replaced.

(6) Your answer packet contains only 4 sheets of graph paper; no further graph paper will be provided. After the experiment you may take the slider away.

(7) If the timer cannot be used, ask the invigilator for help.

\textbf{[Data]}

\begin{center}
Weight of the slider $=5.84\times10^{-2}\unitv{N}$\\[2pt]
The reading of the voltmeter represents the time for which the slider slides between the two sensors; $1\unitv{V}$ corresponds to $0.213\unitv{s}$ (when the potential difference of the battery is $9\unitv{V}$)\\[2pt]
Distance between the two sensors $=0.294\unitv{m}$
\end{center}

\textbf{[Contents of the experiment]}

(1) Using only the equipment provided, investigate how $\xi_{d'}$ depends on the speed $v_b$ of the slider on an inclined track making an angle $\theta$ with the horizontal.

(2) On the answer sheet, state the algebraic equations or relations used to analyse the experimental results and to plot the graphs.

(3) Propose a quantitative model to explain the experimental results. Use the collected data to justify this model.

\subsection*{Problem 5}

When the magnetic slider slides down an aluminium inclined plane from top to bottom with uniform speed $v$, the induced emf (and hence the induced current) is proportional to the speed $v$, so the magnetic damping force may be taken as $F=Kv$, where $K$ is a constant. This magnetic damping force may also be written as $\xi_dW\cos\theta$, where $W$ is the weight of the slider and $\xi_d$ may be regarded as a ``magnetic'' friction coefficient; $\theta$ is the angle between the inclined plane and the horizontal.

When the slider slides down the incline with uniform speed, the resultant of the forces parallel to the incline is zero, so that
\begin{equation*}
W\sin\theta=Kv+\mu W\cos\theta ,
\end{equation*}
where $\mu$ is the coefficient of kinetic friction between the slider and the incline. Dividing both sides of this equation by $W\cos\theta$ gives the equation
\begin{equation*}
\tan\theta=\frac{Kv}{W\cos\theta}+\mu .
\end{equation*}
Plotting a straight-line graph of $\tan\theta$ against $v/\cos\theta$ gives the slope $K/W$ and the intercept $\mu$; hence the magnetic friction coefficient $\xi_d$ can be found. Its relation to the speed $v$ and the angle $\theta$ is
\begin{equation*}
\xi_d=\text{slope}\cdot\frac{v}{\cos\theta} ,
\end{equation*}
where the value of the slope is obtained from the graph.

\begin{figbox}{31.10}
\begin{tikzpicture}[scale=0.9]
\draw[thickline] (0,0) -- (6,0) -- (6,2.6) -- cycle;
\draw[thin] (0.6,0) arc (0:23.4:0.6);
\node at (1.05,0.22) {$\theta$};
\begin{scope}[rotate=23.4]
\draw[thickline] (2.6,0) -- (3.5,0) -- (3.5,0.45) -- (2.6,0.45) -- cycle;
\draw[->,thickline] (3.1,0.55) -- (2.0,0.55) node[left] {$W\sin\theta$};
\draw[->,thickline] (3.1,0.22) -- (4.2,0.22) node[right] {$Kv+\mu W\cos\theta$};
\end{scope}
\end{tikzpicture}
\end{figbox}

[Note] The raw reading of the timer is a voltage $V$ (the voltage read by the digital multimeter). $V$ can be converted into a time $t$ by the relation $t=0.213V$; hence the speed of the slider can be found as $v=0.294/t$, that is, the speed of the slider is $v=1.380/V$ (where $V$ is the voltage across the capacitor).

This problem is a design-and-investigation competition problem; the manipulation, measurement and analysis should include the following.

(1) Measure the voltage across the laminated battery with the voltage range of the digital multimeter, and then calibrate the voltage across the capacitor in the timer:
\begin{equation*}
V_{\text{correct}}=V_{\text{reading}}\cdot(9.00/\text{battery voltage}).
\end{equation*}

(2) The experiment uses angles of inclination $\theta$ in the range $\theta_{\max}>\theta>\theta_{\min}$. Show that within this range the slider moves with uniform speed between the two sensors; this can be shown by recording the same time measured by the timer when the slider starts from different heights. Record these data.

(3) At different angles of inclination $\theta$, let the slider slide down with uniform speed, record the time $t$ measured by the timer, and find the speed $v$ of the slider as it passes between the two integrated Hall sensors. Change the angle $\theta$ and take measurements at at least 5 points; plot a straight-line graph of $\tan\theta$ against $v/\cos\theta$. Multiplying the slope of the line by $v/\cos\theta$ gives $\xi_d$, and the intercept is the coefficient of kinetic friction $\mu$; the result is that $\mu$ should lie in the range $0.2$--$0.4$. A layer of paper is stuck to both the aluminium incline and the surface of the slider; do not touch them with your hands.
